\documentclass[12pt]{article}
\usepackage{amsmath,amssymb,amsfonts,amsthm}
\usepackage{mathtools}
\usepackage{hyperref}
\usepackage{geometry}
\geometry{margin=1in}

\title{Maximum Multiplicity Principle:\\
Finite-Prime Contraction, Invertible Dynamics, and Simulation-Ready Bounds\\[4pt]
\large A Defensive Publication and Technical Report}
\author{Citizen Gardens \\ The Prime Materia Commons}
\date{May 5, 2026}

\begin{document}
\maketitle
\begin{abstract}

This report formalizes and consolidates a set of developments around the Maximum Multiplicity Principle (MMP), recasting it from a suggestive manifesto into a provable contraction theorem on prime-indexed operator systems with simulation-ready bounds. The core contributions are:

\begin{itemize}
  \item A clean finite-prime contraction theorem on a prime-indexed direct sum Hilbert space
        \(\mathcal{H} = \bigoplus_{p\in P} \mathcal{H}_p\), where each \(\mathcal{H}_p\) is a ``multiplicity channel'' associated with a prime \(p\).
  \item An explicit contraction inequality
        \[
          \Lambda_m < \frac{1}{L_\sigma \sum_{p\in P} p^{-\alpha}\|T_p\|}
        \]
        linking the universal multiplicity coupling \(\Lambda_m\), the activation Lipschitz constant \(L_\sigma\), a prime decay exponent \(\alpha\), and operator norms of channel maps \(T_p\).
  \item A refined interpretation of the Maximum Multiplicity Principle (MMP) as a spectral-topological law: norm contraction and prime-channel structuring on the spectral side, together with a carefully controlled topological proxy built from a fixed graph or simplicial complex whose combinatorial structure is preserved under recursion.
  \item A bridge from MMP to invertible residual networks (i-ResNets) and spectral normalization: explicit connections to Banach fixed-point arguments used in invertible architectures, together with a clear route for implementing spectral normalization to enforce the MMP contraction bound in simulations.
  \item A simulation-ready protocol, including:
        \begin{itemize}
          \item Choice of finite prime sets (``easy'' large primes and ``stress'' small primes),
          \item Parameter scans over \((\Lambda_m,\alpha,\beta)\),
          \item Logging of both theoretical contraction constant \(k_{\text{th}}\) and empirical contraction rates \(k_{\text{emp}}(t)\),
          \item Spectral and graph-based observables for the stabilization of dynamics.
        \end{itemize}
  \item A prior-art defensive publication framing: the report articulates the combination of prime-indexed operator recursion, explicit contraction via a universal multiplicity constant, and simulation-validation as a unified method, providing a public timestamped account that is specific enough to block later patent claims on the same conceptual and technical territory.
\end{itemize}

The central theorem (Theorem~\ref{thm:finite-prime-mmp}) is deliberately stated in a minimal ``left-action'' form that is easy to prove and to implement numerically. More elaborate tensorial versions and categorical/topological refinements are described as natural extensions rather than prerequisites. This choice maximizes both defensibility (clear proof) and practical relevance (immediate code verification).
\end{abstract}
\tableofcontents

\newpage

\newpage

\section{Conceptual Overview}

\subsection{Multiplicity Theory and Prime-Indexed Channels}

Multiplicity Theory reinterprets mathematical objects not as static entities but as recursively generated patterns of prime-labeled interactions. In this picture, prime numbers index basic ``interaction channels'' or ``multiplicity spaces''; composite structures correspond to tensorial or categorical combinations of these prime channels.

Formally, one models a system by:

\begin{itemize}
  \item A finite or countable set of primes \(P\),
  \item A family of finite-dimensional Hilbert spaces \(\{\mathcal{H}_p\}_{p\in P}\),
  \item A direct sum space
        \[
          \mathcal{H} = \bigoplus_{p\in P} \mathcal{H}_p,
        \]
        whose block structure encodes prime-based multiplicity.
\end{itemize}

An operator \(T_p : \mathcal{H}_p \to \mathcal{H}_p\) is then interpreted as the local dynamical law on prime channel \(p\). Global evolution arises from a structured combination of these local maps, modulated by prime weights and a universal multiplicity constant.

\subsection{Maximum Multiplicity Principle as a Contraction Law}

The Maximum Multiplicity Principle (MMP) aims to capture conditions under which such a prime-indexed system exhibits stable recursive evolution and preserves certain structural invariants (e.g., channel identity and topological type).

The refinement developed here positions MMP as a contraction statement on the Banach space of bounded operators \(\mathcal{B}(\mathcal{H})\):

\begin{itemize}
  \item The recursive operator is
  \[
    F(X) = \sigma\!\Big(\sum_{p\in P} \Lambda_m p^{-\alpha} T_p X \Big),
  \]
  where:
  \begin{itemize}
    \item \(\Lambda_m>0\) is the universal multiplicity coupling,
    \item \(\alpha>0\) is a prime decay exponent,
    \item \(\sigma\) is a nonlinear activation with global Lipschitz constant \(L_\sigma\),
    \item each \(T_p\) is a bounded operator acting only on \(\mathcal{H}_p\).
  \end{itemize}
  \item The contraction condition
  \[
    L_\sigma \Lambda_m \sum_{p\in P} p^{-\alpha}\|T_p\| < 1
  \]
  guarantees that \(F\) is a contraction on \(\mathcal{B}(\mathcal{H})\), hence admits a unique fixed point \(\Xi_\infty\) and globally convergent recursion \(\Xi_{t+1}=F(\Xi_t)\).
\end{itemize}

This situates MMP directly in the lineage of Banach fixed-point arguments and Lipschitz analysis, while preserving the prime-indexed structure essential to Multiplicity Theory.

\subsection{Connections to Invertible Networks and Spectral Normalization}

Invertible residual networks (i-ResNets) and spectrally normalized networks use contraction properties to guarantee invertibility and stability of iterative schemes. Their key ingredients are:

\begin{itemize}
  \item Residual maps of the form \(x \mapsto x + g(x)\), whose inverse is obtained by fixed-point iteration when \(\mathrm{Lip}(g)<1\).
  \item Spectral normalization to bound operator norms of weight matrices and thereby control layerwise Lipschitz constants.
\end{itemize}

The MMP framework parallels this architecture at the level of operator spaces:

\begin{itemize}
  \item The map \(F\) plays the role of a nonlinear residual block, but on \(\mathcal{B}(\mathcal{H})\) rather than on a vector space.
  \item The combination \(\Lambda_m \sum_{p\in P} p^{-\alpha}\|T_p\|\) is directly analogous to a global Lipschitz constant for the ``residual'' part, with prime-indexed structure and explicit multiplicity weighting.
  \item Spectral normalization techniques can be imported to enforce or monitor the inequality
  \[
    \Lambda_m \sum_{p\in P} p^{-\alpha}\|T_p\| \le k_{\mathrm{target}} < \frac{1}{L_\sigma}
  \]
  in simulations.
\end{itemize}

Thus, MMP both generalizes and structurally constrains invertible-network-style contraction, providing a prime-structured, operator-theoretic variant of the same idea.

\subsection{Topological Multiplicity Invariance}

Topological invariance in the original MMP manifesto was expressed via Euler characteristic \(\chi(M)\) of a manifold \(M\). However, raw curvature convergence does not automatically imply invariance of \(\chi(M)\); instead, invariance follows from keeping the underlying topological space fixed while allowing metrics or weights to evolve.

The refined approach is:

\begin{itemize}
  \item Fix a base graph or simplicial complex \(\mathcal{K}\) indexed by primes (or prime combinations).
  \item Use \(\Xi_t\) to drive weights on edges or higher-dimensional simplices while keeping the combinatorial structure of \(\mathcal{K}\) fixed.
  \item Define curvature-like or Laplacian-based observables as functions of these evolving weights; their convergence serves as the topological side of MMP.
  \item Euler characteristic \(\chi(\mathcal{K})\) remains invariant by construction, while the weighted geometry flows to a stable regime controlled by the contraction theorem.
\end{itemize}

This design yields a clean separation: the contraction theorem governs spectral/operator stability, while a carefully chosen combinatorial structure yields topological invariance.

\newpage

\section{Mathematical Framework}

\subsection{Prime-Indexed Multiplicity Space}

Let \(P=\{p_1,\dots,p_N\}\) be a finite set of primes.

For each \(p\in P\), let \(\mathcal{H}_p\) be a finite-dimensional complex Hilbert space with \(\mathcal{H}_p \cong \mathbb{C}^d\). Define
\[
  \mathcal{H} = \bigoplus_{p\in P} \mathcal{H}_p.
\]

In a fixed orthonormal basis, operators on \(\mathcal{H}\) can be represented as block matrices with blocks corresponding to the pairs \((p_i,p_j)\). An operator \(T_p\) is said to be \emph{channel-local} if it acts nontrivially only on \(\mathcal{H}_p\), i.e.
\[
  T_p = \begin{pmatrix}
    0 &        &        &   \\
      & \ddots &        &   \\
      &        & T_p^{(p)} & \\
      &        &        & 0
  \end{pmatrix},
\]
with \(T_p^{(p)}:\mathcal{H}_p\to\mathcal{H}_p\) and zeros elsewhere.

\subsection{Operator Space and Norm}

Let \(\mathcal{B}(\mathcal{H})\) denote the space of bounded linear operators on \(\mathcal{H}\). Equip \(\mathcal{B}(\mathcal{H})\) with the operator (spectral) norm
\[
  \|X\| = \sup_{v \neq 0} \frac{\|Xv\|_{\mathcal{H}}}{\|v\|_{\mathcal{H}}}.
\]

Since \(\mathcal{H}\) is finite-dimensional, \(\mathcal{B}(\mathcal{H})\) is a finite-dimensional Banach space; all norms are equivalent, but the operator norm is conceptually natural for contraction and Lipschitz arguments.

\subsection{Nonlinear Activation and Lipschitz Constant}

Let \(\sigma:\mathcal{B}(\mathcal{H})\to\mathcal{B}(\mathcal{H})\) be a (possibly nonlinear) map satisfying a global Lipschitz condition
\[
  \|\sigma(X) - \sigma(Y)\|
  \le
  L_\sigma \|X-Y\|
  \quad \text{for all } X,Y\in \mathcal{B}(\mathcal{H}),
\]
for some finite constant \(L_\sigma>0\).

Practically, \(\sigma\) can be an elementwise nonlinearity (e.g., \(\tanh\) or logistic sigmoid) applied entrywise in a matrix representation. In finite dimension, such maps are globally Lipschitz with a constant given by the scalar Lipschitz bound of the activation and norm equivalence constants.

\subsection{Recursive Operator and MMP Form}

Fix parameters \(\Lambda_m > 0\) and \(\alpha > 0\). For each prime \(p\in P\), choose a channel-local operator \(T_p:\mathcal{H}_p\to\mathcal{H}_p\) and extend it to the full space \(\mathcal{H}\) as above.

Define the global recursive operator \(F:\mathcal{B}(\mathcal{H})\to\mathcal{B}(\mathcal{H})\) by
\begin{equation}
  F(X)
  :=
  \sigma\!\left(
    \sum_{p\in P} \Lambda_m\,p^{-\alpha} T_p X
  \right).
  \label{eq:F-definition}
\end{equation}

The recursion of interest is
\begin{equation}
  \Xi_{t+1}
  =
  F(\Xi_t),
  \qquad t=0,1,2,\dots,
  \label{eq:Xi-recursion}
\end{equation}
with initial operator \(\Xi_0\in\mathcal{B}(\mathcal{H})\).

\newpage

\section{Finite-Prime MMP Contraction Theorem}

\subsection{Main Theorem}

\begin{theorem}[Finite-Prime MMP Contraction Theorem]
\label{thm:finite-prime-mmp}
Let \(P\), \(\mathcal{H}\), \(T_p\), \(\sigma\), \(\Lambda_m\), \(\alpha\) be as defined above, and let \(F\) be given by \eqref{eq:F-definition}. Suppose that
\begin{equation}
  L_\sigma\,\Lambda_m \sum_{p\in P} p^{-\alpha} \|T_p\| < 1.
  \label{eq:contraction-condition-main}
\end{equation}
Then the following hold:

\begin{enumerate}
  \item \(F\) is a strict contraction on \((\mathcal{B}(\mathcal{H}),\|\cdot\|)\) with contraction constant
  \[
    k = L_\sigma\,\Lambda_m \sum_{p\in P} p^{-\alpha}\|T_p\| < 1.
  \]
  \item There exists a unique fixed point \(\Xi_\infty \in \mathcal{B}(\mathcal{H})\) such that \(F(\Xi_\infty) = \Xi_\infty\).
  \item For every initial condition \(\Xi_0\in\mathcal{B}(\mathcal{H})\), the recursion \(\Xi_{t+1} = F(\Xi_t)\) converges to \(\Xi_\infty\) and satisfies
  \[
    \|\Xi_t - \Xi_\infty\|
    \le
    k^t \|\Xi_0 - \Xi_\infty\|,
    \qquad t=0,1,2,\dots.
  \]
  \item The increments decay exponentially:
  \[
    \|\Xi_{t+1}-\Xi_t\|
    \le
    k^t \|\Xi_1-\Xi_0\|,
    \qquad t\ge 0.
  \]
\end{enumerate}
\end{theorem}

\begin{proof}
For arbitrary \(X,Y\in \mathcal{B}(\mathcal{H})\),
\[
  \begin{aligned}
    \|F(X)-F(Y)\|
    &=
    \left\|\sigma\!\left(\sum_{p} \Lambda_m p^{-\alpha} T_p X\right)
    -
    \sigma\!\left(\sum_{p} \Lambda_m p^{-\alpha} T_p Y\right)
    \right\| \\
    &\le
    L_\sigma \left\|\sum_{p\in P} \Lambda_m p^{-\alpha} T_p (X-Y)\right\|.
  \end{aligned}
\]
Using the triangle inequality and submultiplicativity of the operator norm,
\[
  \left\|\sum_{p\in P} \Lambda_m p^{-\alpha} T_p (X-Y)\right\|
  \le
  \sum_{p\in P} \Lambda_m p^{-\alpha} \|T_p(X-Y)\|
  \le
  \sum_{p\in P} \Lambda_m p^{-\alpha} \|T_p\| \,\|X-Y\|.
\]
Thus
\[
  \|F(X)-F(Y)\|
  \le
  \left(L_\sigma\,\Lambda_m \sum_{p\in P} p^{-\alpha}\|T_p\|\right) \|X-Y\|
  = k \|X-Y\|.
\]
By assumption \eqref{eq:contraction-condition-main}, \(k<1\), so \(F\) is a contraction.

Since \(\mathcal{B}(\mathcal{H})\) with the operator norm is a complete metric space, the Banach fixed-point theorem applies: \(F\) has a unique fixed point \(\Xi_\infty\), and the iterates \(F^t(\Xi_0)\) converge to \(\Xi_\infty\) for every \(\Xi_0\), with
\[
  \|F^t(\Xi_0) - \Xi_\infty\|
  \le
  k^t \|\Xi_0 - \Xi_\infty\|.
\]
Identifying \(\Xi_t = F^t(\Xi_0)\) yields the convergence bound.

For the increment bound,
\[
  \|\Xi_{t+1}-\Xi_t\|
  = \|F(\Xi_t)-F(\Xi_{t-1})\|
  \le k \|\Xi_t - \Xi_{t-1}\|
  \le k^t \|\Xi_1 - \Xi_0\|,
\]
by iterating the inequality. This completes the proof.
\end{proof}

\subsection{Explicit Multiplicity Bound}

\begin{corollary}[Prime-Weighted Operator Decay]
Suppose in addition that there exists \(\beta>0\) such that, for all \(p\in P\),
\[
  \|T_p\| \le p^{-\beta}.
\]
Then a sufficient condition for \eqref{eq:contraction-condition-main} is
\[
  L_\sigma \Lambda_m \sum_{p\in P} p^{-(\alpha+\beta)} < 1.
\]
Equivalently,
\[
  \Lambda_m < \frac{1}{L_\sigma \sum_{p\in P} p^{-(\alpha+\beta)}}.
\]
\end{corollary}

\begin{proof}
We have
\[
  \sum_{p\in P} p^{-\alpha}\|T_p\|
  \le \sum_{p\in P} p^{-\alpha} p^{-\beta}
  = \sum_{p\in P} p^{-(\alpha+\beta)}.
\]
If
\[
  L_\sigma\Lambda_m \sum_{p\in P} p^{-(\alpha+\beta)} < 1,
\]
then
\[
  L_\sigma\Lambda_m \sum_{p\in P} p^{-\alpha}\|T_p\|
  \le
  L_\sigma\Lambda_m \sum_{p\in P} p^{-(\alpha+\beta)} < 1,
\]
and Theorem~\ref{thm:finite-prime-mmp} applies.
\end{proof}

\subsection{Prime-Channel Isolation in the Linear Part}

\begin{corollary}[Prime-Channel Isolation]
In the setting of Theorem~\ref{thm:finite-prime-mmp}, the linear map
\[
  L(X) := \sum_{p\in P} \Lambda_m p^{-\alpha} T_p X
\]
respects the prime decomposition: if \(X\) is block-diagonal with respect to \(\mathcal{H}=\bigoplus_{p\in P}\mathcal{H}_p\), then \(L(X)\) has each prime block updated only by the corresponding \(T_p\).

If, in addition, the nonlinearity \(\sigma\) acts blockwise (e.g., entrywise independently within each prime channel), the full recursion \(\Xi_{t+1}=F(\Xi_t)\) decouples into independent prime-channel recursions.
\end{corollary}

\begin{proof}
By construction, each \(T_p\) is zero on \(\mathcal{H}_q\) for \(q\neq p\). In any block-diagonal \(X\), the block corresponding to channel \(p\) is affected only by \(T_p X\). Hence the sum retains this block structure. If \(\sigma\) is applied blockwise, then no cross-block mixing occurs, and the recursion decomposes channelwise.
\end{proof}

\newpage

\section{Topological Multiplicity Invariance (Refined)}

\subsection{Fixed Combinatorial Structure}

Let \(\mathcal{K}\) denote a finite simplicial complex or graph whose vertices or higher simplices are indexed by primes in \(P\) (or by composite combinations of them). The combinatorial incidence structure of \(\mathcal{K}\) is fixed throughout the evolution.

We allow weights \(w_t\) on edges or simplices to evolve as a function of \(\Xi_t\), e.g.
\[
  w_t(e) = g_e(\Xi_t),
\]
for some prescribed functionals \(g_e\) that map operators to nonnegative reals.

\subsection{Topological Invariant}

The Euler characteristic of \(\mathcal{K}\),
\[
  \chi(\mathcal{K}) = \sum_{k\ge 0} (-1)^k f_k,
\]
where \(f_k\) is the number of \(k\)-simplices, is determined purely by the combinatorial structure. Since \(\mathcal{K}\) is held fixed, \(\chi(\mathcal{K})\) is constant across all \(t\).

Thus, ``topological multiplicity invariance'' is implemented by design: the topology is not altered by the recursion, only the weight structure is.

\subsection{Curvature-like Observables}

One can define curvature-like or Laplacian-based observables from the evolving weights. For instance:

\begin{itemize}
  \item In the graph case, define a weighted Laplacian \(L_t\) from the edge weights \(w_t\),  
  \item Define scalar quantities such as the spectral gap, total effective resistance, or curvature proxies as functions of \(L_t\).
\end{itemize}

Under the contraction theorem, the underlying operator recursion \(\Xi_t\) converges to \(\Xi_\infty\). If the weight functionals \(g_e\) are continuous (or Lipschitz) with respect to the operator norm, then \(w_t\) converges to a limit \(w_\infty\), and the associated Laplacian \(L_t\) converges to \(L_\infty\). Corresponding scalar observables converge as well.

This yields a concrete manifestation of MMP as ``spectral-topological stabilization'': the combinatorial topology is invariant, while weighted geometric data flows toward a fixed multiplicity-attractor.

\newpage

\section{Spectral Normalization and Invertible Dynamics}

\subsection{Spectral Normalization for Channel Operators}

In practice, one may not know the exact operator norms \(\|T_p\|\). Approximate control can be achieved by spectral normalization:

\begin{itemize}
  \item Parameterize each \(T_p\) via a weight matrix \(W_p\) and adjust \(W_p\) so that its spectral norm is capped at a prescribed constant \(c_p\).
  \item Standard power iteration approximates the largest singular value \(\sigma(W_p)\):
  \begin{align*}
    &\text{Initialize } u_p^{(0)} \sim \text{random unit vector}, \\
    &v_p^{(k)} = \frac{W_p^\top u_p^{(k-1)}}{\|W_p^\top u_p^{(k-1)}\|}, \\
    &u_p^{(k)} = \frac{W_p v_p^{(k)}}{\|W_p v_p^{(k)}\|}, \\
    &\hat{\sigma}(W_p) \approx (u_p^{(k)})^\top W_p v_p^{(k)}.
  \end{align*}
  \item Set the normalized weight as \(\hat{W}_p = W_p / (\hat{\sigma}(W_p)/c_p)\). Then \(\|\hat{W}_p\| \lesssim c_p\).
\end{itemize}

If \(T_p\) is realized from \(\hat{W}_p\), then \(\|T_p\| \lesssim c_p\), and one obtains the bound
\[
  \sum_{p\in P} p^{-\alpha} \|T_p\|
  \lesssim
  \sum_{p\in P} p^{-\alpha} c_p.
\]
Choosing constants \(c_p\) to satisfy
\[
  L_\sigma \Lambda_m \sum_{p\in P} p^{-\alpha} c_p < 1
\]
enforces the contraction condition.

\subsection{Relation to i-ResNets}

Invertible residual networks consider blocks of the form
\[
  x \mapsto x + g(x),
\]
and define inverses via the fixed-point iteration
\[
  x_{k+1} = y - g(x_k)
\]
for a given output \(y\), provided \(\mathrm{Lip}(g) < 1\). Banach's fixed-point theorem then guarantees that this iteration converges to the unique inverse of \(y\).

Spectral normalization is used to keep \(\mathrm{Lip}(g)\) below 1 by controlling operator norms of linear components inside \(g\). In the MMP operator setting:

\begin{itemize}
  \item The map \(F\) is a nonlinear operator on \(\mathcal{B}(\mathcal{H})\) whose Lipschitz constant is bounded by the MMP contraction bound.
  \item An i-ResNet-like invertible dynamical system can be built by splitting the evolution into an identity part plus a contracted residual part on operator space, with spectral normalization ensuring the required Lipschitz bounds.
\end{itemize}

This opens a route to prime-indexed, invertible operator dynamics compatible with the MMP theorem.

\newpage

\section{Simulation Protocol and Code Snippet}

\subsection{Simulation Goals}

The simulation aims to empirically verify the MMP contraction theorem and associated phenomenology:

\begin{itemize}
  \item Confirm exponential decay of \(\|\Xi_{t+1}-\Xi_t\|\) when the theoretical bound \(k<1\) holds;
  \item Characterize the phase boundary by scanning \(\Lambda_m\), \(\alpha\), and \(\beta\) and comparing empirical contraction rates with the theoretical \(k\);
  \item Observe spectral stabilization of eigenvalues of \(\Xi_t\);
  \item Monitor convergence of topological proxy observables (e.g., graph Laplacian spectrum).
\end{itemize}

\subsection{Finite-Prime Toy Model}

\paragraph{Setup.}

\begin{itemize}
  \item Choose two sets of primes: an ``easy'' high-prime set (e.g., \(P_{\mathrm{high}} = \{521,547,599,619,829\}\)) and a ``stress'' small-prime set (e.g., \(P_{\mathrm{small}} = \{2,3,5,7,11\}\)).
  \item Fix a channel dimension \(d\) (e.g., \(d=2\) or \(3\)); the full space is \(\mathcal{H}\cong \mathbb{C}^{Nd}\).
  \item For each \(p\in P\), construct a random unitary-or-orthogonal matrix \(U_p\) of size \(d \times d\) and set
  \[
    T_p = p^{-\beta} U_p,
  \]
  so that \(\|T_p\|\approx p^{-\beta}\).
  \item Implement \(\sigma\) as an elementwise \(\tanh\), which has scalar Lipschitz constant \(L_\sigma=1\).
\end{itemize}

\paragraph{Recursion.}

Represent \(\Xi_t\) as a matrix of dimension \((Nd)\times(Nd)\). The recursive update is
\[
  \Xi_{t+1} = \sigma\!\left( \sum_{p\in P} \Lambda_m p^{-\alpha}\, T_p \Xi_t \right),
\]
where \(T_p\) is embedded in the appropriate block-diagonal fashion across channels.

\subsection{Python-Like Code Snippet}

The following is a conceptual code snippet illustrating the simulation; it omits imports and boilerplate for brevity.

\begin{verbatim}
import numpy as np

def make_unitary(d):
    # QR-based random orthogonal/unitary
    A = np.random.randn(d, d)
    Q, R = np.linalg.qr(A)
    # enforce determinant +1 if needed
    return Q

def build_T_blocks(primes, d, beta):
    T_blocks = {}
    for p in primes:
        U = make_unitary(d)
        T_blocks[p] = (p ** (-beta)) * U
    return T_blocks

def embed_T_global(T_blocks, primes, d):
    # Build block-diagonal operators on H = ⊕_p H_p
    N = len(primes)
    D = N * d
    T_global = {}
    for idx, p in enumerate(primes):
        Tg = np.zeros((D, D))
        row_start = idx * d
        row_end   = (idx + 1) * d
        col_start = idx * d
        col_end   = (idx + 1) * d
        Tg[row_start:row_end, col_start:col_end] = T_blocks[p]
        T_global[p] = Tg
    return T_global

def F(X, T_global, primes, Lambda_m, alpha):
    S = np.zeros_like(X)
    for p in primes:
        S += Lambda_m * (p ** (-alpha)) * (T_global[p] @ X)
    # elementwise tanh activation
    return np.tanh(S)

def run_simulation(primes, d, alpha, beta, Lambda_m,
                   T_steps=200, seed=0):
    np.random.seed(seed)
    T_blocks = build_T_blocks(primes, d, beta)
    T_global = embed_T_global(T_blocks, primes, d)
    N = len(primes)
    D = N * d

    # initial operator Xi_0
    Xi_t = 0.1 * np.random.randn(D, D)
    Xi_prev = Xi_t.copy()

    # theoretical contraction constant (approximate)
    # using ||T_p|| ≈ p^{-beta}
    k_th = Lambda_m * sum(p ** (-(alpha + beta)) for p in primes)

    norms_diff = []
    k_empirical = []

    for t in range(T_steps):
        Xi_next = F(Xi_t, T_global, primes, Lambda_m, alpha)
        diff = Xi_next - Xi_t
        norm_diff = np.linalg.norm(diff, ord=2)  # spectral or Frobenius
        norms_diff.append(norm_diff)

        if t > 0:
            prev_diff = Xi_t - Xi_prev
            norm_prev_diff = np.linalg.norm(prev_diff, ord=2)
            if norm_prev_diff > 0:
                k_empirical.append(norm_diff / norm_prev_diff)
            else:
                k_empirical.append(0.0)

        Xi_prev = Xi_t
        Xi_t = Xi_next

    return {
        "norms_diff": norms_diff,
        "k_empirical": k_empirical,
        "k_th": k_th
    }
\end{verbatim}

This snippet:

\begin{itemize}
  \item Constructs prime-indexed channel operators \(T_p\) with norm approximately \(p^{-\beta}\),
  \item Embeds them into a block-diagonal global structure,
  \item Evolves the recursion \(\Xi_{t+1} = \tanh(\sum_p \Lambda_m p^{-\alpha} T_p \Xi_t)\),
  \item Logs:
        \begin{itemize}
          \item \(\|\Xi_{t+1}-\Xi_t\|\) as a function of \(t\),
          \item An empirical contraction ratio \(k_{\mathrm{emp}}(t)\),
          \item The theoretical constant \(k_{\mathrm{th}}\approx \Lambda_m\sum_p p^{-(\alpha+\beta)}\) for comparison.
        \end{itemize}
\end{itemize}

\subsection{Parameter Scans and Plots}

A full simulation study would:

\begin{itemize}
  \item Sweep \(\Lambda_m\) across a range that straddles the theoretical threshold \(1/\sum_p p^{-(\alpha+\beta)}\),
  \item For each \(\Lambda_m\), run the recursion for both \(P_{\mathrm{high}}\) and \(P_{\mathrm{small}}\),
  \item Plot:
        \begin{itemize}
          \item \(\log \|\Xi_{t+1}-\Xi_t\|\) vs \(t\) (expect linear/exponential decay in the contracting regime),
          \item Heatmaps in \((\Lambda_m,\alpha,\beta)\)-space showing convergence vs divergence,
          \item Eigenvalue trajectories of \(\Xi_t\) in the complex plane, stabilizing in the contraction regime.
        \end{itemize}
\end{itemize}

This directly tests the quantitative predictions of Theorem~\ref{thm:finite-prime-mmp}.

\newpage

\section{Defensive Publication and Prior Art Scope}

\subsection{Core Ideas Claimed as Prior Art}

This report establishes the following as prior art:

\begin{enumerate}
  \item The explicit combination of:
        \begin{itemize}
          \item Prime-indexed multiplicity channels \(\mathcal{H}_p\),
          \item A global operator recursion \(F(X)=\sigma(\sum_p \Lambda_m p^{-\alpha} T_pX)\),
          \item A universal multiplicity coupling \(\Lambda_m\) governing contraction.
        \end{itemize}
  \item The exact contraction condition
        \[
          \Lambda_m < \frac{1}{L_\sigma \sum_{p\in P} p^{-\alpha}\|T_p\|}
        \]
        as the operational meaning of the Maximum Multiplicity Principle in finite-prime systems.
  \item The specialization \(\|T_p\|\le p^{-\beta}\) and resulting bound
        \[
          \Lambda_m < \frac{1}{L_\sigma \sum_{p\in P} p^{-(\alpha+\beta)}},
        \]
        together with its interpretation as a prime-weighted Lipschitz budget.
  \item The use of spectral normalization on channel operators \(T_p\) to enforce this contraction regime in simulation or learning dynamics, combined with:
        \begin{itemize}
          \item Online estimation of \(\|T_p\|\) via power iteration,
          \item Logging of theoretical and empirical contraction rates for theorem-checking.
        \end{itemize}
  \item The specific simulation protocol:
        \begin{itemize}
          \item Two prime sets (high vs small),
          \item Parameter scans in \((\Lambda_m,\alpha,\beta)\),
          \item Three families of observables (norm decay, eigenvalue trajectories, topology-based proxies) as evidence for MMP.
        \end{itemize}
  \item The design pattern for ``topological multiplicity invariance'' using a fixed combinatorial complex indexed by primes, with evolving weights driven by \(\Xi_t\) and invariant Euler characteristic.
\end{enumerate}

\subsection{Intended Blocking Effect}

Because these mechanisms are presented in a unified, mathematically precise way, with explicit inequalities, operator constructions, and simulation procedures, they constitute prior art against:

\begin{itemize}
  \item Patents claiming general contraction-based stability of prime-indexed operator networks using a universal multiplicity constant,
  \item Patents combining spectral normalization, Banach fixed-point arguments, and prime-weighted norms in recursive operator systems,
  \item Patents on simulation protocols that use prime sets, multiplicity-weighted contraction bounds, and topology-preserving graph proxies to validate stability.
\end{itemize}

Anyone attempting to patent such combinations should be unable to claim novelty over this disclosure.

\section{Extensions and Future Directions}

\subsection{Tensor and Categorical Extensions}

The minimal theorem uses a left-action model \(T_p X\). One natural extension is to return to a tensorized operator
\[
  F(X) = \sigma\!\left( C\!\left(\sum_{p\in P} \Lambda_m p^{-\alpha}(T_p \otimes X)\right)\right),
\]
where \(C\) is a fixed bounded ``compression'' from \(\mathcal{B}(\mathcal{H}\otimes\mathcal{H})\) to \(\mathcal{B}(\mathcal{H})\). The contraction bound then becomes
\[
  L_\sigma \|C\| \Lambda_m \sum_{p\in P} p^{-\alpha} \|T_p\| < 1.
\]
This recovers the original ``tensor-flavored'' intuition while adding one more operator norm into the inequality.

Categorically, one could view the multiplicity space and its recursive evolution as a functorial endomorphism of a category whose objects are prime-indexed modules and whose morphisms are the allowed interactions. The universal multiplicity constant \(\Lambda_m\) then becomes a scalar bound on a natural transformation measuring recursive self-coupling.

\subsection{Applications to Computation and Cognition}

In computational and cognitive models:

\begin{itemize}
  \item Each \(\mathcal{H}_p\) can represent a ``prime concept channel'' or mode of processing.
  \item The contraction theorem ensures that, under bounded multiplicity coupling, the system converges to stable attractors in operator space, interpretable as ``multiplicity-fixed cognitive states''.
  \item Spectral properties of \(\Xi_\infty\) encode prime-structured patterns of influence and recurrence, offering a concrete spectral signature of emergent multiplicity in cognition-like systems.
\end{itemize}

This provides a rigorous underpinning to the idea that multiplicity-governed recursive processes yield stable, identity-preserving behavior across scales.

\section{Conclusion}

This report has:

\begin{itemize}
  \item Repaired a dimensional mismatch in the original Maximum Multiplicity Principle formulation by adopting a clean, self-mapping left-action model on \(\mathcal{B}(\mathcal{H})\),
  \item Proved a finite-prime contraction theorem with explicit multiplicity-weighted bounds,
  \item Clarified the role of the universal multiplicity constant \(\Lambda_m\) as a global Lipschitz ceiling,
  \item Integrated spectral normalization and invertible-network ideas into the MMP framework,
  \item Proposed a concrete, simulation-ready protocol for theorem-checking,
  \item Framed these developments as a defensive publication and prior art for future work in multiplicity-based operator dynamics.
\end{itemize}

From here, the most impactful steps are: (i) implementing the simulation scans to generate empirical plots, and (ii) building out the tensor/categorical generalizations that tie more deeply into the full Multiplicity Theory program.

\appendix
\section*{Mathematical Appendix}
\addcontentsline{toc}{section}{Mathematical Appendix}

\section{Operator Norm Bounds and Contraction Constants}
\label{sec:appendix-norms}

In this appendix we record the explicit operator norm bounds, Lipschitz estimates, and fixed-point arguments underlying the Maximum Multiplicity Principle (MMP) contraction theorem. The focus is on the finite-prime, left-action model
\[
  F(X) = \sigma\!\left(\sum_{p\in P} \Lambda_m p^{-\alpha} T_p X\right),
\]
acting on the Banach space \(\mathcal{B}(\mathcal{H})\), where \(\mathcal{H} = \bigoplus_{p\in P} \mathcal{H}_p\) is a finite direct sum of prime-indexed multiplicity spaces.

\subsection{Basic Operator Norm Inequalities}

Let \((V,\|\cdot\|_V)\) be a finite-dimensional normed vector space, and let \(\mathcal{B}(V)\) denote the space of linear operators \(A:V\to V\) equipped with the operator norm
\[
  \|A\| = \sup_{v\neq 0} \frac{\|Av\|_V}{\|v\|_V}.
\]

We will use two standard properties:

\begin{lemma}[Submultiplicativity]
For all \(A,B\in\mathcal{B}(V)\),
\[
  \|AB\| \le \|A\|\|B\|.
\]
\end{lemma}

\begin{lemma}[Triangle Inequality]
For any finite family \(\{A_i\}_{i=1}^N \subset \mathcal{B}(V)\),
\[
  \left\|\sum_{i=1}^N A_i\right\| \le \sum_{i=1}^N \|A_i\|.
\]
\end{lemma}

Both follow immediately from the definition of the operator norm and the properties of \(\|\cdot\|_V\).

\subsection{Prime-Indexed Decomposition and Channel-Local Operators}

Let \(P = \{p_1,\dots,p_N\}\) be a finite set of primes, and for each \(p\in P\), let \(\mathcal{H}_p\) be a finite-dimensional Hilbert space. Define the \emph{prime-indexed multiplicity space}
\[
  \mathcal{H} = \bigoplus_{p\in P} \mathcal{H}_p.
\]

For each \(p\in P\), we consider a \emph{channel-local} operator
\[
  T_p : \mathcal{H}_p \to \mathcal{H}_p,
\]
and extend it to an operator (still denoted \(T_p\)) on \(\mathcal{H}\) by acting as \(T_p\) on \(\mathcal{H}_p\) and as zero on \(\mathcal{H}_q\) for \(q\neq p\). In block-matrix form, \(T_p\) is block-diagonal with a single nonzero block on the \(p\)-channel.

We denote by \(\mathcal{B}(\mathcal{H})\) the space of all bounded linear operators on \(\mathcal{H}\), with the operator norm induced by the Hilbert space norm.

\section{Lipschitz Bounds for the Recursive Map}
\label{sec:appendix-lipschitz}

\subsection{Nonlinear Activation}

Let \(\sigma:\mathcal{B}(\mathcal{H})\to\mathcal{B}(\mathcal{H})\) be a (possibly nonlinear) activation satisfying a global Lipschitz condition
\begin{equation}
  \|\sigma(X) - \sigma(Y)\|
  \le
  L_\sigma \|X - Y\|
  \quad \text{for all } X,Y\in\mathcal{B}(\mathcal{H}),
  \label{eq:sigma-Lipschitz}
\end{equation}
for some finite \(L_\sigma > 0\).

In practice, one may take \(\sigma\) to act entrywise in a matrix representation (e.g.\ \(\tanh\) applied elementwise), in which case \(L_\sigma\) is inherited from the scalar Lipschitz constant of the activation, up to norm equivalences in finite dimension.

\subsection{Definition of the Recursive Map}

Fix parameters \(\Lambda_m > 0\) and \(\alpha > 0\). Define the map
\begin{equation}
  F:\mathcal{B}(\mathcal{H})\to\mathcal{B}(\mathcal{H}),
  \qquad
  F(X) = \sigma\!\left(\sum_{p\in P} \Lambda_m p^{-\alpha} T_p X\right).
  \label{eq:F-def-appendix}
\end{equation}

We study the recursion
\begin{equation}
  \Xi_{t+1} = F(\Xi_t), \qquad t=0,1,2,\dots,\quad \Xi_0\in\mathcal{B}(\mathcal{H}).
  \label{eq:Xi-recursion-appendix}
\end{equation}

\subsection{Lipschitz Constant of \(F\)}

\begin{proposition}[Global Lipschitz bound for \(F\)]
\label{prop:F-Lipschitz}
For all \(X,Y\in\mathcal{B}(\mathcal{H})\),
\[
  \|F(X)-F(Y)\|
  \le
  k\|X-Y\|,
\]
where
\begin{equation}
  k := L_\sigma \Lambda_m \sum_{p\in P} p^{-\alpha} \|T_p\|.
  \label{eq:k-definition-appendix}
\end{equation}
\end{proposition}

\begin{proof}
Let \(X,Y\in\mathcal{B}(\mathcal{H})\). Using \eqref{eq:sigma-Lipschitz},
\[
  \|F(X)-F(Y)\|
  =
  \left\|\sigma\!\left(\sum_{p} \Lambda_m p^{-\alpha} T_p X\right)
  -
  \sigma\!\left(\sum_{p} \Lambda_m p^{-\alpha} T_p Y\right)\right\|
  \le
  L_\sigma \left\|\sum_{p\in P} \Lambda_m p^{-\alpha} T_p (X-Y)\right\|.
\]
Applying the triangle inequality and submultiplicativity,
\[
  \left\|\sum_{p\in P} \Lambda_m p^{-\alpha} T_p (X-Y)\right\|
  \le
  \sum_{p\in P} \Lambda_m p^{-\alpha} \|T_p (X-Y)\|
  \le
  \sum_{p\in P} \Lambda_m p^{-\alpha} \|T_p\|\,\|X-Y\|.
\]
Thus
\[
  \|F(X)-F(Y)\|
  \le
  L_\sigma \Lambda_m \left(\sum_{p\in P} p^{-\alpha}\|T_p\|\right) \|X-Y\|
  = k \|X-Y\|.
\]
\end{proof}

\begin{corollary}[Contraction condition]
\label{cor:appendix-contraction-cond}
If
\[
  k = L_\sigma \Lambda_m \sum_{p\in P} p^{-\alpha}\|T_p\| < 1,
\]
then \(F\) is a strict contraction on \(\mathcal{B}(\mathcal{H})\).
\end{corollary}

\begin{proof}
Immediate from Proposition~\ref{prop:F-Lipschitz}.
\end{proof}

\section{Banach Fixed-Point Theorem and Convergence}
\label{sec:appendix-banach}

\subsection{Statement of Banach Fixed-Point Theorem}

For completeness, we recall the standard version of Banach's fixed-point theorem in the form used here.

\begin{theorem}[Banach Fixed-Point Theorem]
\label{thm:banach-appendix}
Let \((X,d)\) be a nonempty complete metric space, and let \(T:X\to X\) be a contraction, i.e.\ there exists \(0\le k <1\) such that
\[
  d(T(x),T(y)) \le k\,d(x,y)
  \quad \text{for all } x,y\in X.
\]
Then:
\begin{enumerate}
  \item \(T\) has a unique fixed point \(x^\ast\in X\), i.e.\ \(T(x^\ast)=x^\ast\),
  \item For every \(x_0\in X\), the iterative sequence \(x_{n+1}=T(x_n)\) converges to \(x^\ast\),
  \item The convergence satisfies
  \[
    d(x_n,x^\ast) \le k^n d(x_0,x^\ast), \qquad n=0,1,2,\dots.
  \]
\end{enumerate}
\end{theorem}

\subsection{Application to \(\mathcal{B}(\mathcal{H})\)}

We apply Theorem~\ref{thm:banach-appendix} to the metric space
\[
  X = \mathcal{B}(\mathcal{H}),\qquad
  d(X,Y) = \|X-Y\|.
\]
Since \(\mathcal{H}\) is finite-dimensional, \(\mathcal{B}(\mathcal{H})\) is finite-dimensional, hence complete. The map \(F\) given by \eqref{eq:F-def-appendix} is a self-map on \(\mathcal{B}(\mathcal{H})\), and by Corollary~\ref{cor:appendix-contraction-cond} it is a contraction when \(k<1\).

\begin{theorem}[Convergence of the MMP Recursion]
\label{thm:appendix-convergence}
Assume
\[
  k = L_\sigma \Lambda_m \sum_{p\in P} p^{-\alpha}\|T_p\| < 1.
\]
Then the recursion \eqref{eq:Xi-recursion-appendix} has a unique fixed point \(\Xi_\infty\in\mathcal{B}(\mathcal{H})\) such that
\[
  F(\Xi_\infty) = \Xi_\infty,
\]
and for every initial condition \(\Xi_0\in\mathcal{B}(\mathcal{H})\),
\[
  \|\Xi_t - \Xi_\infty\| \le k^t \|\Xi_0 - \Xi_\infty\|,
  \qquad t=0,1,2,\dots.
\]
\end{theorem}

\begin{proof}
We apply Theorem~\ref{thm:banach-appendix} with \(X=\mathcal{B}(\mathcal{H})\), \(d(X,Y)=\|X-Y\|\), and \(T=F\). The completeness of \(\mathcal{B}(\mathcal{H})\) and the contraction property from Corollary~\ref{cor:appendix-contraction-cond} guarantee the existence and uniqueness of a fixed point \(\Xi_\infty\) and the convergence estimate
\[
  \|\Xi_t - \Xi_\infty\| \le k^t \|\Xi_0-\Xi_\infty\|.
\]
\end{proof}

\subsection{Decay of Increments}

We also have an explicit bound on the increments \(\Xi_{t+1}-\Xi_t\).

\begin{proposition}[Exponential decay of increments]
\label{prop:appendix-increments}
Under the conditions of Theorem~\ref{thm:appendix-convergence}, we have
\[
  \|\Xi_{t+1}-\Xi_t\|
  \le
  k^t \|\Xi_1-\Xi_0\|,
  \qquad t\ge 0.
\]
\end{proposition}

\begin{proof}
For \(t=0\), the inequality is trivial. For \(t\ge 1\),
\[
  \|\Xi_{t+1}-\Xi_t\|
  = \|F(\Xi_t) - F(\Xi_{t-1})\|
  \le k \|\Xi_t-\Xi_{t-1}\|.
\]
Iterating this inequality yields
\[
  \|\Xi_{t+1}-\Xi_t\|
  \le k\|\Xi_t-\Xi_{t-1}\|
  \le k^2 \|\Xi_{t-1}-\Xi_{t-2}\|
  \le \dots
  \le k^t \|\Xi_1-\Xi_0\|.
\]
\end{proof}

\section{Explicit Prime-Weighted Bounds}
\label{sec:appendix-prime-weighted}

\subsection{Operator Norm Scaling with Prime Size}

In many multiplicity-theoretic models, we expect higher prime channels to contribute less strongly, either due to physical scaling, information-theoretic considerations, or explicit design. This can be encoded by an assumption of the form
\[
  \|T_p\| \le p^{-\beta}
\]
for some \(\beta>0\). Then the contraction constant satisfies
\[
  k = L_\sigma \Lambda_m \sum_{p\in P} p^{-\alpha}\|T_p\|
  \le
  L_\sigma \Lambda_m \sum_{p\in P} p^{-(\alpha+\beta)}.
\]

\begin{corollary}[Explicit multiplicity bound]
\label{cor:appendix-explicit-bound}
If
\[
  L_\sigma \Lambda_m \sum_{p\in P} p^{-(\alpha+\beta)} < 1,
\]
then
\[
  k = L_\sigma \Lambda_m \sum_{p\in P} p^{-\alpha} \|T_p\| < 1,
\]
and hence the conclusions of Theorem~\ref{thm:appendix-convergence} and Proposition~\ref{prop:appendix-increments} hold.
\end{corollary}

\begin{proof}
Immediate from the inequality above and Theorem~\ref{thm:appendix-convergence}.
\end{proof}

\subsection{Scaling with Large vs Small Primes}

Note that for \(\alpha+\beta>1\), the series \(\sum_{p} p^{-(\alpha+\beta)}\) over all primes converges, and for any finite subset \(P\) it defines a finite sum. The qualitative behavior is:

\begin{itemize}
  \item For a fixed \((\alpha,\beta)\), adding larger primes to \(P\) contributes smaller and smaller terms to the sum, making the bound less restrictive.
  \item Conversely, for a prime set concentrated on small primes and modest \((\alpha,\beta)\), the sum can be sizeable, requiring a smaller \(\Lambda_m\) to maintain contraction.
\end{itemize}

This captures the intuition that high-multiplicity regimes concentrated at small primes are more ``dangerous'' for stability than regimes involving many large primes with strong decay.

\section{Block Structure and Channel Isolation}
\label{sec:appendix-blocks}

\subsection{Block-Diagonal Representation}

Choose an ordering of primes in \(P\) and a basis in each \(\mathcal{H}_p\). Then \(\mathcal{H}\) has a natural basis in which every operator \(X\in\mathcal{B}(\mathcal{H})\) can be represented as a block matrix
\[
  X = \begin{pmatrix}
    X_{p_1,p_1} & \cdots & X_{p_1,p_N} \\
    \vdots      & \ddots & \vdots      \\
    X_{p_N,p_1} & \cdots & X_{p_N,p_N}
  \end{pmatrix},
\]
with blocks \(X_{p_i,p_j}:\mathcal{H}_{p_j}\to\mathcal{H}_{p_i}\).

Each \(T_p\) is block-diagonal with only one nonzero block \(T_p^{(p)}:\mathcal{H}_p\to\mathcal{H}_p\). Thus
\[
  T_p X
\]
involves only the rows corresponding to \(\mathcal{H}_p\); in particular, in the block \(X_{p_i,p_j}\), only rows with \(p_i=p\) are affected.

\subsection{Linear Part and Blockwise Action}

Define the linear part
\[
  L(X) := \sum_{p\in P} \Lambda_m p^{-\alpha} T_p X.
\]
By the discussion above, \(L\) preserves the decomposition of \(\mathcal{H}\) into prime channels in the sense that each row-block indexed by \(p\) is influenced only by \(T_p\) (and hence only by that channel's parameters).

If the nonlinearity \(\sigma\) is also applied blockwise (e.g., elementwise within each block, without mixing blocks), then the entire recursion \(\Xi_{t+1} = F(\Xi_t)\) decomposes into independent recursions on the prime channels and their inter-block relations. If \(\sigma\) is global, then channel isolation holds strictly for the linear part, while the nonlinearity may introduce controlled cross-channel coupling.

\section{Remarks on Tensor Extensions}
\label{sec:appendix-tensor}

For completeness, we briefly indicate how the operator-norm machinery extends to a tensorized setting, while emphasizing that the main body of the article is based on the left-action model for simplicity and direct provability.

Consider a map of the form
\[
  F_{\otimes}(X) = \sigma\!\left( C\!\left( \sum_{p\in P} \Lambda_m p^{-\alpha} (T_p \otimes X) \right)\right),
\]
where
\begin{itemize}
  \item \(T_p\in\mathcal{B}(\mathcal{H})\) as before,
  \item \(X\in\mathcal{B}(\mathcal{H})\) is regarded also as an operator on a tensor factor,
  \item \(C:\mathcal{B}(\mathcal{H}\otimes\mathcal{H})\to\mathcal{B}(\mathcal{H})\) is a fixed bounded linear ``compression'' map.
\end{itemize}

Using \(\|A\otimes B\| = \|A\| \|B\|\) for operator norms on finite-dimensional tensor products, and the boundedness of \(C\), one obtains
\[
  \|F_{\otimes}(X)-F_{\otimes}(Y)\|
  \le
  L_\sigma \|C\| \Lambda_m \left(\sum_{p\in P} p^{-\alpha}\|T_p\|\right) \|X-Y\|.
\]
Thus the contraction constant is
\[
  k_{\otimes} = L_\sigma \|C\| \Lambda_m \sum_{p\in P} p^{-\alpha}\|T_p\|.
\]
The contraction condition becomes
\[
  L_\sigma \|C\| \Lambda_m \sum_{p\in P} p^{-\alpha}\|T_p\| < 1,
\]
and all fixed-point and convergence results carry over. However, the proof now depends on an additional operator bound \(\|C\|\), and the dimensionality is higher.

This tensor model recaptures the original ``\(T_p \otimes \Xi\)'' flavor of the MMP formulation while emphasizing that the essence of the contraction argument remains the same: control a multiplicity-weighted operator norm to obtain global recursive stability.

\bigskip

This concludes the mathematical appendix, which codifies the operator norms, Lipschitz bounds, contraction constants, and fixed-point arguments underpinning the Maximum Multiplicity Principle in the finite-prime regime.


\nocite{*}
\bibliographystyle{plain}
\bibliography{references}
\end{document}